The GlobalMinimum library of Jonas Mockus (1989) implements some of the
earliest practical Bayesian global-optimization methods. We modernize
this Fortran 66/77 codebase along both standard routes: wrapping the
original compiled routines behind R and Python packages, and
translating the core methods to Rust. First, we release packages in
which every optimizer runs the genuine 1989 algorithms, with callback
objectives, deterministic seeding, evaluation traces, and validation
that retrofits memory- and I/O-safety onto COMMON-block Fortran.
Second, we dissect a translation-fidelity incident: the Rust port ran
up to 12x slower than the wrapped Fortran---not from foreign-function
overhead, but from two dropped pruning rules whose restoration recovers
bitwise trajectory agreement and speed parity. The same audit surfaced
transposed test-function matrices and an uninitialized read in the
upstream code itself. Third, we benchmark the 1989 methods against
modern optimizers in both ecosystems (78,000 runs, externally validated
on BBOB/COCO, bitwise-replicated across languages): the coordinate
Wiener-model search EXKOR statistically ties SHGO, dual annealing, and
GenSA on classical test functions but drops to mid-field on BBOB's
non-separable landscapes, and a compound method built from the
library's own components trades aggregate solve rate for markedly
fewer total failures. All code and data are available for replication.
